% Clase 24 --- La onda electromagnetica plana propuesta (§4.3).
% E en y, B en z, ambos funcion de (x,t) y EN FASE. Perspectiva oblicua: el eje
% z (hacia el lector) se dibuja bajando a la izquierda.
% Nota: el sub-dibujo del par de vectores se repite, por eso va con providecommand.
\providecommand{\campos}[1]{%
  \pgfmathsetmacro{\amp}{sin(deg(#1))}%
  \draw[figvec] (#1,0) -- (#1,{1.35*\amp});
  \draw[figvecres] (#1,0) -- ({#1-0.45*\amp},{-0.32*\amp});}
\begin{center}
\begin{tikzpicture}[scale=1.02]

  % ---- ejes ----
  \draw[figaxis] (-0.6,0) -- (7.6,0) node[figlblaux, anchor=north] {$x$};
  \draw[figaxis] (0,-1.7) -- (0,1.85) node[figlblaux, anchor=east] {$y$};
  \draw[figaxis] (0,0) -- (-1.15,-0.82) node[figlblaux, anchor=east] {$z$};

  % ---- envolventes ----
  \draw[figline, smooth, samples=120, domain=0:7.2]
    plot (\x, {1.35*sin(deg(\x))});
  \draw[figred, line width=0.9pt, smooth, samples=120, domain=0:7.2]
    plot ({\x-0.45*sin(deg(\x))}, {-0.32*sin(deg(\x))});

  % ---- vectores en posiciones equiespaciadas ----
  \foreach \xx in {0.4,0.8,...,7.2}{\campos{\xx}}

  % ---- rótulos ----
  \node[figlbl, figblue, anchor=south west] at (1.35,1.3) {$\mathbf{E} \;(\hat y)$};
  \node[figlbl, figred, anchor=north east] at (0.95,-0.42) {$\mathbf{B} \;(\hat z)$};

  \draw[figvecres, line width=1.4pt] (7.6,1.35) -- (8.55,1.35);
  \node[figlbl, anchor=south] at (8.05,1.42) {$c$};

\end{tikzpicture}\\[0.5em]
{\small\itshape Ésta es la familia de soluciones que se propone ---no la más
general: se verifica que satisface las ecuaciones de Maxwell y con eso alcanza
para calcular su velocidad. Las hipótesis están todas dibujadas: la onda es
\textbf{transversal} (ambos campos perpendiculares a la dirección de
propagación $\hat x$ y perpendiculares entre sí), es \textbf{plana} (los campos
no dependen de $y$ ni de $z$: todo el plano perpendicular a $\hat x$ oscila
junto, y por eso los vectores dibujados sobre el eje representan lo que pasa en
todo ese plano) y los dos campos están \textbf{en fase}, es decir que donde uno
se anula el otro también. Conviene mirar los ceros: ahí no hay campo de ninguna
especie, y sin embargo un instante después vuelve a haberlo, sostenido por lo
que ocurre en las zonas vecinas. Eso es lo notable de esta onda: no hay medio
que la sostenga, cada campo variable genera al otro y así se propaga sola por el
vacío ---exactamente lo que hace la luz del Sol durante todo el viaje hasta la
Tierra.}
\end{center}
